\subsection{Coordinates and Matrix Representations}
Throughout this subsection, all vector spaces are over the same fixed field $\mathbb{F}$.

\begin{definition} \label{def:coordinate-vector-and-coordinate-map} \hlt{Coordinate Vector}\\
Let $n\geq1$, let $\beta=(\mathbf{v}_1,\ldots,\mathbf{v}_n)$ be a basis of a vector space $\mathcal{V}$, and let $\mathbf{v}\in\mathcal{V}$. By Proposition \ref{prop:basis-unique-coordinate-representation}, $(\exists!\,(c_1,\ldots,c_n)\in\mathbb{F}^n)\,\mathbf{v}=\sum\limits_{j=1}^{n}c_j\mathbf{v}_j$.\\
For this tuple, the \hlt{$\beta$-coordinate vector} of $\mathbf{v}$ is $[\mathbf{v}]_{\beta}\defeq(c_1,\ldots,c_n)^t\in\mathbb{F}^n$.\\
The function $[\cdot]_{\beta}:\mathcal{V}\to\mathbb{F}^n$ is the \hlt{$\beta$-coordinate map}. The order of the list $\beta$ fixes the order of the coordinates.
\end{definition}

\begin{proposition} \label{prop:coordinate-map-isomorphism} \hlt{Coordinate Maps Are Isomorphisms}\\
Let $n\geq1$, and let $\beta=(\mathbf{v}_1,\ldots,\mathbf{v}_n)$ be a basis of a vector space $\mathcal{V}$. The coordinate map $[\cdot]_{\beta}:\mathcal{V}\to\mathbb{F}^n$ has the following properties:
\begin{enumerate}[label=(\roman*)]
\setlength{\itemsep}{0pt}
\item It is linear.
\item It is injective.
\item It is surjective and hence an isomorphism.
\end{enumerate}
\begin{proof}
\begin{enumerate}[label=(\roman*)]
\setlength{\itemsep}{0pt}
\item Let $\mathbf{u},\mathbf{v}\in\mathcal{V}$ and $c,d\in\mathbb{F}$, and write $[\mathbf{u}]_{\beta}=(a_1,\ldots,a_n)^t$ and $[\mathbf{v}]_{\beta}=(b_1,\ldots,b_n)^t$.\\
Then $[c\mathbf{u}+d\mathbf{v}]_{\beta}=(ca_1+db_1,\ldots,ca_n+db_n)^t=c[\mathbf{u}]_{\beta}+d[\mathbf{v}]_{\beta}$, so $[\cdot]_{\beta}$ is linear.
\item If $[\mathbf{u}]_{\beta}=[\mathbf{v}]_{\beta}$, then $\mathbf{u}$ and $\mathbf{v}$ have the same coefficients in $\beta$, so $\mathbf{u}=\mathbf{v}$ by Proposition \ref{prop:basis-unique-coordinate-representation}. Thus $[\cdot]_{\beta}$ is injective.
\item For $(c_1,\ldots,c_n)^t\in\mathbb{F}^n$, set $\mathbf{v}\defeq\sum\limits_{j=1}^{n}c_j\mathbf{v}_j$. Then $[\mathbf{v}]_{\beta}=(c_1,\ldots,c_n)^t$, so $[\cdot]_{\beta}$ is surjective.\\
It is therefore an isomorphism by Proposition \ref{prop:linear-map-invertible-iff-bijective}.
\end{enumerate}
\hlt{Structural Commentary}: The three checks compare vectors with their unique coefficient tuples.\\
\hlt{Key Insight}: An ordered basis identifies the vector space with its coordinate space without losing linear structure.
\end{proof}
\end{proposition}

\begin{definition} \label{def:matrix-representation-of-linear-map} \hlt{Matrix Representation}\\
Let $n,m\geq1$, let $\beta=(\mathbf{v}_1,\ldots,\mathbf{v}_n)$ be a basis of $\mathcal{V}$, and let $\gamma=(\mathbf{w}_1,\ldots,\mathbf{w}_m)$ be a basis of $\mathcal{W}$.\\
For $T\in\LL(\mathcal{V},\mathcal{W})$, the \hlt{matrix of $T$ from $\beta$-coordinates to $\gamma$-coordinates} is $[T]_{\gamma\leftarrow\beta}\defeq\bigl[[T(\mathbf{v}_1)]_{\gamma}\ \cdots\ [T(\mathbf{v}_n)]_{\gamma}\bigr]\in\operatorname{Mat}(m,n,\mathbb{F})$.\\
Thus its $j$th column is the uniquely determined coordinate vector $[T(\mathbf{v}_j)]_{\gamma}$.
\end{definition}

\begin{proposition} \label{prop:matrix-representation-respects-linear-combinations} \hlt{Matrices Respect Linear Combinations of Maps}\\
Let $\beta=(\mathbf{v}_1,\ldots,\mathbf{v}_n)$ be a basis of $\mathcal{V}$, and let $\gamma=(\mathbf{w}_1,\ldots,\mathbf{w}_m)$ be a basis of $\mathcal{W}$.\\
If $S,T\in\LL(\mathcal{V},\mathcal{W})$ and $a\in\mathbb{F}$, then $[aS+T]_{\gamma\leftarrow\beta}=a[S]_{\gamma\leftarrow\beta}+[T]_{\gamma\leftarrow\beta}$.
\begin{proof}
By Proposition \ref{prop:linear-maps-form-vector-space}, $aS+T\in\LL(\mathcal{V},\mathcal{W})$. For $j\in\{1,\ldots,n\}$, $[(aS+T)\mathbf{v}_j]_{\gamma}=[aS\mathbf{v}_j+T\mathbf{v}_j]_{\gamma}=a[S\mathbf{v}_j]_{\gamma}+[T\mathbf{v}_j]_{\gamma}$ by Proposition \ref{prop:coordinate-map-isomorphism}(i). These are the respective $j$th columns of the two matrices, so the matrices are equal.\\
\hlt{Structural Commentary}: The proof compares the coordinate images of each domain basis vector.\\
\hlt{Key Insight}: Taking a matrix representation preserves the vector-space operations on linear maps.
\end{proof}
\end{proposition}

\begin{proposition} \label{prop:matrix-representation-action} \hlt{Matrix Representation Acts on Coordinates}\\
Let $\beta=(\mathbf{v}_1,\ldots,\mathbf{v}_n)$ be a basis of $\mathcal{V}$, let $\gamma=(\mathbf{w}_1,\ldots,\mathbf{w}_m)$ be a basis of $\mathcal{W}$, and let $T\in\LL(\mathcal{V},\mathcal{W})$. Then $(\forall\,\mathbf{v}\in\mathcal{V})\,[T(\mathbf{v})]_{\gamma}=[T]_{\gamma\leftarrow\beta}[\mathbf{v}]_{\beta}$.
\begin{proof}
Let $\mathbf{v}\in\mathcal{V}$, and write $[\mathbf{v}]_{\beta}=(a_1,\ldots,a_n)^t$, so $\mathbf{v}=\sum\limits_{j=1}^{n}a_j\mathbf{v}_j$. By linearity of $T$ and Proposition \ref{prop:coordinate-map-isomorphism}(i), $[T(\mathbf{v})]_{\gamma}=\sum\limits_{j=1}^{n}a_j[T(\mathbf{v}_j)]_{\gamma}$.\\
The summands are the columns of $[T]_{\gamma\leftarrow\beta}$, and this sum equals $[T]_{\gamma\leftarrow\beta}[\mathbf{v}]_{\beta}$ by Proposition \ref{prop:matrix-vector-product-column-combination}.\\
\hlt{Structural Commentary}: The proof expands a vector in the domain basis and tracks that expansion through $T$.\\
\hlt{Key Insight}: Applying a linear map becomes matrix multiplication after both bases are fixed.
\end{proof}
\end{proposition}
